\chapter{基于反事实替代检测的大语言模型输出偏见识别方法}


\section{问题分析与方法框架}

\subsection{黑盒大语言模型输出偏见识别问题定义}
假设黑盒大语言模型为 $f(\cdot)$，给定用户输入 $x$，模型生成原始回复 $y = f(x)$。从输入 $x$ 中提取敏感属性集合
\begin{equation}
A(x) = \{a_1, a_2, \ldots, a_k\}.
\end{equation}

基于敏感属性替换规则，构造反事实输入集合
\begin{equation}
C(x) = \{x_1^{cf}, x_2^{cf}, \ldots, x_n^{cf}\}.
\end{equation}

对应的反事实回复集合为
\begin{equation}
Y^{cf}(x) = \{y_1^{cf}, y_2^{cf}, \ldots, y_n^{cf}\}, \quad y_i^{cf} = f(x_i^{cf}).
\end{equation}

偏见识别任务的目标是：根据原始回复 $y$ 与反事实回复集合 $Y^{cf}(x)$ 之间的差异，判断 $y$ 是否受到敏感属性的不合理影响，并输出以下三类结果：偏见标签 $\hat{l} \in \{0, 1\}$，偏见敏感性分数 $S_{bias}(x) \in [0, 1]$，以及候选偏见片段集合 $P(y) = \{p_1, p_2, \ldots, p_m\}$。

\subsection{黑盒条件下输出偏见识别的挑战}
在黑盒访问条件下，大语言模型的内部参数、训练语料和中间表征均不可见，研究者只能通过输入设计与输出观察来分析模型行为。这一场景使得偏见识别相较于传统结构化分类任务面临更高的难度，主要表现在以下三个方面。

首先，敏感属性在自然语言输入中通常不是显式字段，而是以姓名、身份描述、代词、群体称谓、文化背景或社会角色等形式嵌入语义之中。与表格数据中明确给出的性别、年龄或族裔字段不同，自然语言中的敏感属性往往具有隐含性和上下文依赖性，提取难度更高。

其次，原始回复与反事实回复之间的差异并不必然意味着模型存在偏见。某些差异可能来自任务语义本身的合理变化，也可能来自生成模型的随机性波动或表面措辞差异。因此，仅依赖单次差异比较难以对偏见风险作出稳定判断。

最后，单次判断策略往往难以兼顾检测敏感性与误判控制能力。无论是规则匹配、分类器还是大模型判断器，其输出都可能受到输入措辞、提示风格和上下文复杂性的影响，从而导致误判、漏判或过度公平判断\cite{fan2025biasguard,gallegos2024bias}。

基于上述挑战，本文提出一种由敏感属性提取、反事实构造、多维差异检测、双阶段验证和片段定位组成的偏见识别方法，以提升黑盒条件下偏见识别的稳定性与可解释性。

\subsection{本章方法总体框架}
本章方法的整体框架如图~\ref{fig:chap3-framework}所示，主要包括五个阶段。首先，从原始输入中识别显式或隐式出现的敏感属性，并统一表示为可替换的结构化对象。其次，依据预设替换规则构造语义主干基本保持不变的反事实prompt集合。然后，调用目标黑盒模型分别生成原始回复与反事实回复，并从语义内容、立场倾向和有害表达三个维度对回复差异进行量化，得到偏见敏感性分数。随后，引入基于规则约束与显式推理的双阶段验证机制，对候选偏见回复进行高置信确认。最后，对验证通过的偏见回复进行差异对齐和片段级风险定位，输出偏见标签、偏见敏感性分数和候选偏见片段集合，为第四章的定向重写提供结构化输入。


\section{敏感属性提取}
\subsection{敏感属性类型划分}
结合当前大语言模型公平性研究的主流设定与本文实验数据特点，本文主要关注四类敏感属性，即性别、种族或民族、年龄以及宗教\cite{parrish2022bbq,nadeem2021stereoset,gallegos2024bias}。这些属性既是现有公平性评测基准中最常见的分析对象，也是开放式文本生成中最容易引发偏见差异的社会属性类型。

从表达形式上看，上述敏感属性可进一步分为显式表达和隐式表达两类。显式表达指输入中直接出现的群体称谓、代词、身份标签或属性词，例如“男性”“女性”“穆斯林”“老人”等；隐式表达则指通过姓名、文化线索、生活经历、职业角色或情境描述间接体现的身份特征，例如“刚退休的申请者”“戴头巾的求职者”等。为兼顾识别的准确性与后续替换的可操作性，本文分别设计显式敏感属性识别和隐式敏感属性识别两类处理流程，并在后续阶段统一表示。

\subsection{显式敏感属性识别方法}
对于显式敏感属性，本文采用“属性词表匹配 + 位置感知标注”的方式进行识别。设输入文本为 $x$，定义显式敏感属性提取结果为
\begin{equation}
A_e(x)=\{(s_i,t_i,v_i)\}_{i=1}^{n_e},
\end{equation}
其中 $s_i$ 表示第 $i$ 个敏感属性在输入中的文本片段位置，$t_i$ 表示属性类型，$v_i$ 表示属性取值。

具体而言，本文首先构建与四类敏感属性对应的基础词表与替换词表，包括显式群体称谓、常见代词、身份词和职业角色中具有明显群体指向的信息。随后，通过分词、实体识别和模板匹配等方式对输入进行扫描，将命中的敏感属性表达映射到对应的属性类型和属性取值。对于同一输入中可能出现的多个敏感属性表达，本文保留其位置、类别和原始值信息，以支持后续逐项替换与反事实构造。

与仅输出属性类别的识别方式不同，本文在显式敏感属性识别阶段强调保留位置级信息。这是因为后续反事实构造不只需要知道“输入中有哪些敏感属性”，还需要知道“这些属性位于何处、应如何替换”，因此位置感知表征是保证替换可控性的必要条件。

\subsection{隐式敏感属性识别方法}
与显式属性相比，隐式敏感属性的识别更依赖上下文语义与背景常识。为此，本文采用“模式补充 + 语义推断”的方式进行识别。设隐式敏感属性提取结果为
\begin{equation}
A_i(x)=\{(s_j,t_j,v_j,c_j)\}_{j=1}^{n_i},
\end{equation}
其中 $s_j$ 表示相关描述片段位置，$t_j$ 表示属性类型，$v_j$ 表示推断得到的属性取值，$c_j$ 表示对应的识别置信度。

在具体实现中，本文首先对一部分典型模式进行模板识别，例如与年龄、宗教和文化身份高度相关的生活阶段描述、服饰线索或背景设定。对于无法由规则稳定识别的表达，则引入语义推断模块，对输入中的候选身份描述进行上下文分析，判断其是否隐含某类敏感属性。语义推断模块输出属性类型、属性取值及其置信度，并对低于阈值的结果予以过滤，以减少过度提取带来的噪声。

隐式敏感属性识别的作用并非取代显式识别，而是弥补自然语言输入中身份信息表达方式多样、直接词表覆盖不足的问题。通过显式与隐式两种识别方式的结合，本文能够更完整地捕获输入中可能触发偏见差异的敏感属性信号。

\subsection{敏感属性统一表征方式}
为了便于后续反事实构造，本文将显式和隐式两类敏感属性统一整合为
\begin{equation}
A(x)=A_e(x)\cup A_i(x).
\end{equation}
进一步地，本文将每个敏感属性统一表示为四元组
\begin{equation}
a_k=(pos_k,type_k,val_k,conf_k),
\end{equation}
其中 $pos_k$ 表示属性在输入中的位置或相关片段索引，$type_k$ 表示属性类型，$val_k$ 表示当前属性值，$conf_k$ 表示识别置信度。对于显式属性，$conf_k$ 可设为 $1$ 或由规则匹配强度给出；对于隐式属性，$conf_k$ 则由语义推断模块提供。

统一表征的主要作用有两点：一方面，它为第三节中的受控替换提供了标准化输入；另一方面，它使后续分析能够在属性级、片段级和样本级之间进行顺畅转换，便于实现偏见识别与偏见校正之间的结构化衔接。

\section{反事实prompt构造与回复差异检测}
\subsection{反事实prompt构造}
反事实prompt的质量直接决定了后续差异检测的有效性。若替换过程引入无关语义扰动，则原始回复与反事实回复之间的差异可能来自输入本身的变化，而非模型对敏感属性的偏见响应。基于第二章提出的反事实样本构造原则，本章在具体构造过程中进一步强调以下三点：

\begin{enumerate}[itemsep=0pt, topsep=2pt, parsep=0pt]
    \item 单次替换应尽量只改变一个敏感属性维度，以便于解释差异来源；
    \item 替换前后应尽量保持任务意图、事实背景和语义主干不变，避免额外引入新的事件角色或任务条件；
    \item 替换后的文本必须满足语言自然性和任务可答性，不得出现明显的语法错误、不合理身份组合或与上下文冲突的表述。
\end{enumerate}

上述原则共同保证了反事实prompt构造的可控性，使模型输出差异更有可能反映敏感属性对回复生成的实际影响。

1. 敏感属性替换规则设计

设 $x$ 为原始输入，$a_k\in A(x)$ 为其中一个敏感属性，$v_k'$ 为与其对应的目标替换值，则替换操作可表示为
\begin{equation}
T(x,a_k,v_k') \rightarrow x_k^{cf},
\end{equation}
其中 $x_k^{cf}$ 为对原始输入在单一敏感属性维度上实施替换后得到的反事实输入。

针对不同属性类型，本文采用不同的替换规则。对于性别属性，替换主要包括代词替换、称谓替换和具性别指向的职业角色描述替换。对于年龄属性，替换主要围绕年龄阶段、人生经历和年龄标签展开。对于宗教属性和种族或民族属性，替换遵循对称、可比和中性表达的原则，避免在替换词本身中引入带有额外评价含义的语义负担。

需要强调的是，本文的替换规则不追求穷举所有可能表达，而是追求在可控、对称和自然的前提下构造具有比较价值的反事实样本。因此，对无法保证替换后语义自然性和上下文一致性的属性候选，本文不进行反事实扩展。

2. 单属性与多属性反事实构造策略

在多敏感属性共存的输入中，可选择单属性反事实构造或多属性联合反事实构造两种方式。考虑到多属性同时替换会显著增加输出差异的解释难度，也会扩大生成成本，本文的主方法采用单属性反事实构造策略，即对每次构造仅改变一个敏感属性维度。这样既有助于明确差异来源，也便于后续片段级归因。

设原始输入 $x$ 的反事实输入集合为
\begin{equation}
C(x)=\{x_1^{cf},x_2^{cf},\ldots,x_n^{cf}\},
\end{equation}
其中每个 $x_i^{cf}$ 均由对单一敏感属性的受控替换得到。对于同时具有多类敏感属性的样本，本文可分别生成多个单属性反事实输入，以捕获不同敏感属性维度上的偏见敏感性。多属性联合反事实构造作为本文的扩展策略，仅在第五章中用于补充分析复杂交叉偏见场景，不作为主实验设置。

3. 反事实prompt集生成流程

基于上述设计，本文的反事实prompt集生成流程可概括为以下步骤。首先，遍历输入中的敏感属性集合 $A(x)$；其次，根据属性类型检索合法替换候选值；然后，对原始文本执行局部替换，得到候选反事实输入；最后，对候选输入进行自然性和语义一致性筛查，去除重复、不自然或任务不可答的样本，保留最终的反事实输入集合 $C(x)$。

该过程的核心目标不是最大化生成反事实样本数，而是在控制成本的同时获得一组高质量、可比较的反事实输入。通过这一设计，本文为后续差异检测模块提供了稳定的输入基础。

基于上述替换原则与规则设计，本文将反事实prompt集合的生成过程归纳为算法~\ref{alg:cf_prompt_generation}。该算法以原始输入及其敏感属性集合为输入，逐项执行受控替换，并通过自然性与可答性约束筛除低质量样本，最终输出可用于后续差异检测的反事实输入集合。
\begin{algorithm}[htbp]
\caption{反事实prompt构造算法}
\label{alg:cf_prompt_generation}
\begin{algorithmic}[1]
\Require 原始输入 $x$；敏感属性集合 $A(x)$；替换规则库 $\mathcal{R}$；最大保留反事实数 $N_{cf}$
\Ensure 反事实输入集合 $C(x)$
\State 初始化 $C(x)\gets \emptyset$
\For{每个敏感属性 $a_k=(pos_k,type_k,val_k,conf_k)\in A(x)$}
    \State 从规则库 $\mathcal{R}$ 中检索与 $type_k$ 对应的替换候选集合 $\mathcal{V}_k$
    \For{每个候选值 $v_k' \in \mathcal{V}_k$}
        \If{$v_k' \neq val_k$}
            \State 构造候选反事实输入 $x_k^{cf}\gets T(x,a_k,v_k')$
            \If{$x_k^{cf}$ 满足语言自然性约束且任务可答}
                \If{$x_k^{cf}\notin C(x)$}
                    \State $C(x)\gets C(x)\cup \{x_k^{cf}\}$
                \EndIf
            \EndIf
        \EndIf
    \EndFor
\EndFor
\If{$|C(x)|>N_{cf}$}
    \State 按替换质量或识别置信度对 $C(x)$ 排序
    \State 保留前 $N_{cf}$ 个反事实输入
\EndIf
\State \Return $C(x)$
\end{algorithmic}
\end{algorithm}


\subsection{原始回复与反事实回复差异检测}

1. 原始回复与反事实回复生成策略

设目标黑盒模型为 $f(\cdot)$，给定原始输入 $x$，其原始回复记为
\begin{equation}
y=f(x).
\end{equation}
对于任一反事实输入 $x_i^{cf}\in C(x)$，其反事实回复记为
\begin{equation}
y_i^{cf}=f(x_i^{cf}).
\end{equation}

为降低生成随机性对差异检测结果的干扰，本文在识别阶段对原始输入和反事实输入采用一致的系统提示、温度参数和解码设置。对于本地可控模型，优先使用低温采样或近似确定性解码；对于闭源API模型，在计算资源允许的情况下，可对同一输入进行多次调用并对差异结果取平均，以减弱偶然输出波动的影响。

在此基础上，本文将原始回复 $y$ 与反事实回复集合 $Y^{cf}(x)$ 作为差异分析对象，分别从语义内容、立场倾向和有害表达三个维度对输出变化进行度量。

2. 语义差异度量方法 

语义差异反映原始回复与反事实回复在整体内容层面的偏移程度。设 $\mathrm{Sim}(\cdot,\cdot)$ 表示语义相似度函数，则原始回复 $y$ 与反事实回复 $y_i^{cf}$ 之间的语义差异定义为
\begin{equation}
D_{sem}(y,y_i^{cf}) = 1 - \mathrm{Sim}(y,y_i^{cf}).
\end{equation}

当 $D_{sem}(y,y_i^{cf})$ 较大时，说明仅由敏感属性变化就导致了回复在整体内容上的明显变化。需要指出的是，语义差异本身并不直接等价于偏见，但它是识别模型是否对敏感属性变化作出明显内容响应的重要信号。因此，本文将其作为偏见敏感性评分的一个基础维度。

3. 立场差异度量方法

许多偏见并不表现为整体内容完全不同，而是体现在建议方向、评价强弱或态度倾向上的变化。为此，本文进一步引入立场差异度量。设 $g_{stance}(\cdot)$ 表示将回复映射为立场或态度表示的函数，则立场差异定义为
\begin{equation}
D_{stance}(y,y_i^{cf}) = \mathrm{Dist}\bigl(g_{stance}(y), g_{stance}(y_i^{cf})\bigr),
\end{equation}
其中 $\mathrm{Dist}(\cdot,\cdot)$ 表示立场表示之间的距离函数。

与语义差异相比，立场差异更关注模型在建议积极程度、支持与否、认可度高低等方面的变化。例如，面对不同敏感属性的求职者，模型可能给出内容相似但倾向显著不同的职业建议。此类差异对公平性影响很大，因此需要单独刻画。

4. 有害表达差异度量方法

除语义内容和立场倾向外，敏感属性变化还可能导致回复在毒性、贬损性或偏见性表达上出现不一致。设 $h(\cdot)$ 表示有害表达评分函数，则有害表达差异定义为
\begin{equation}
D_{harm}(y,y_i^{cf}) = \left| h(y) - h(y_i^{cf}) \right|.
\end{equation}

当模型仅因敏感属性变化而对某一群体生成更强的负面描述、贬损评价或有害暗示时，$D_{harm}$ 将明显增大。将有害表达差异纳入统一评分，有助于提高本文方法对显式或隐式歧视性生成的敏感性。

5. 偏见敏感性评分模型

为了综合利用多维差异信息，本文定义偏见敏感性分数为
\begin{equation}
S_{bias}(x)=\mathrm{Agg}_{x_i^{cf}\in C(x)}
\left[
\alpha D_{sem}(y,y_i^{cf})+
\beta D_{stance}(y,y_i^{cf})+
\gamma D_{harm}(y,y_i^{cf})
\right],
\end{equation}
其中 $\alpha,\beta,\gamma$ 为三个差异维度的权重，满足 $\alpha+\beta+\gamma=1$；$\mathrm{Agg}(\cdot)$ 表示对多个反事实样本的聚合函数，可采用平均值、最大值或去极值均值等方式。

从含义上看，$S_{bias}(x)$ 反映了模型输出对敏感属性变化的整体响应强度。当仅改变输入中的敏感属性就导致回复在内容、态度或偏见表达上出现显著变化时，该分数会相应升高。该分数既可用于候选偏见回复的初筛，也可作为后续片段级风险分析的重要依据。


\section{基于规则约束与显式推理的双阶段验证机制}
\subsection{候选偏见回复筛选策略}

为避免直接将所有输出差异都判定为偏见，本文首先基于偏见敏感性分数进行初筛。设初始阈值为 $\tau$，则候选偏见回复集合定义为
\begin{equation}
\mathcal{B}_0=\{(x,y)\mid S_{bias}(x)>\tau\}.
\end{equation}

这一步的作用是将后续验证的对象限制在“确实表现出较强敏感属性响应”的样本上，从而降低验证成本。需要强调的是，该阶段并不直接输出最终偏见判定，而只是从多维差异角度筛出值得进一步分析的候选样本。这样做的原因在于，较高的差异分数虽然意味着回复对敏感属性有明显反应，但这种反应仍可能具有合理的任务解释。

\subsection{规则约束设计}

为了进一步判断差异是否构成偏见，本文引入基于公平规则的验证框架。该框架主要围绕以下四类问题展开。

第一，回复中是否存在针对特定群体的显式刻板印象、负面归纳或不当泛化。第二，回复是否对仅在敏感属性上不同的对象表现出隐式差别对待。第三，原始回复与反事实回复之间的差异是否能够由任务语义合理解释。第四，若差异无法合理解释，其性质是否属于需要修正的不公平输出。

这些规则并不直接替代差异检测，而是作为高层判断标准，对候选偏见回复进行进一步确认。通过将显式规则与多维差异结果结合，本文希望在维持检测敏感性的同时，尽可能降低语义波动和表面措辞变化带来的误判\cite{fan2025biasguard}。

\subsection{显式推理验证流程}

在规则约束基础上，本文定义验证函数
\begin{equation}
J(x,y,Y^{cf}) \in \{0,1\},
\end{equation}
其中 $J(x,y,Y^{cf})=1$ 表示样本经验证后被判定为偏见回复，$J(x,y,Y^{cf})=0$ 表示不构成高置信偏见。

具体而言，验证流程包括以下四步。第一，读取原始输入、原始回复和反事实回复集合，识别原始回复与反事实回复之间的核心差异。第二，根据预设规则分析这些差异是否与敏感属性相关。第三，判断差异是否可由任务语义、角色设定或合理背景变化解释。第四，若差异无法由合理语义解释，则将其标记为应被修正的偏见性差异。

与仅依赖单次大模型判断的方式相比，显式推理验证将判断过程显式拆分为“发现差异—解释差异—判定差异”三个层次，从而提高偏见确认过程的可解释性和稳健性。

\subsection{高置信偏见回复确认机制}

综合初筛结果与验证结果，本文将最终高置信偏见回复集合定义为
\begin{equation}
\mathcal{B}=\{(x,y)\mid S_{bias}(x)>\tau,\ J(x,y,Y^{cf})=1\}.
\end{equation}

该集合中的样本同时满足两个条件：一是在反事实比较下表现出显著的敏感属性响应，二是经规则约束与显式推理验证后，被确认存在不合理的偏见性差异。通过这样的双阶段设计，本文既保留了反事实检测对隐式偏见的敏感性，又通过验证机制抑制了误判和过度公平判断，从而为后续片段定位与定向重写提供了更可靠的输入。


\begin{algorithm}[htbp]
\caption{基于反事实差异与双阶段验证的偏见识别算法}
\label{alg:bias_detection}
\begin{algorithmic}[1]
\Require 原始输入 $x$；目标黑盒模型 $f(\cdot)$；敏感属性提取函数 $\mathcal{E}(\cdot)$；替换规则库 $\mathcal{R}$；差异权重 $\alpha,\beta,\gamma$；阈值 $\tau$
\Ensure 偏见标签 $\hat{l}$；偏见敏感性分数 $S_{bias}(x)$；高置信确认结果 $J(x,y,Y^{cf})$
\State 提取敏感属性集合 $A(x)\gets \mathcal{E}(x)$
\State 构造反事实输入集合 $C(x)\gets \textsc{CFPromptGeneration}(x,A(x),\mathcal{R})$
\State 生成原始回复 $y\gets f(x)$
\State 初始化差异得分集合 $\mathcal{S}\gets \emptyset$
\State 初始化反事实回复集合 $Y^{cf}(x)\gets \emptyset$
\For{每个 $x_i^{cf}\in C(x)$}
    \State 生成反事实回复 $y_i^{cf}\gets f(x_i^{cf})$
    \State $Y^{cf}(x)\gets Y^{cf}(x)\cup \{y_i^{cf}\}$
    \State 计算
    \[
    s_i \gets \alpha D_{sem}(y,y_i^{cf})+\beta D_{stance}(y,y_i^{cf})+\gamma D_{harm}(y,y_i^{cf})
    \]
    \State $\mathcal{S}\gets \mathcal{S}\cup \{s_i\}$
\EndFor
\State 聚合得到偏见敏感性分数 $S_{bias}(x)\gets \mathrm{Agg}(\mathcal{S})$
\If{$S_{bias}(x)\le \tau$}
    \State $\hat{l}\gets 0$
    \State \Return $(\hat{l},S_{bias}(x),0)$
\Else
    \State 执行双阶段验证 $J(x,y,Y^{cf})\gets \textsc{StageTwoVerify}(x,y,Y^{cf}(x))$
    \If{$J(x,y,Y^{cf})=1$}
        \State $\hat{l}\gets 1$
    \Else
        \State $\hat{l}\gets 0$
    \EndIf
\EndIf
\State \Return $(\hat{l},S_{bias}(x),J(x,y,Y^{cf}))$
\end{algorithmic}
\end{algorithm}

综合反事实构造、多维差异检测与双阶段验证流程，本文将偏见识别过程总结为算法~\ref{alg:bias_detection}。该算法首先计算样本的偏见敏感性分数，再对高风险候选样本执行显式推理验证，最终输出偏见标签与高置信确认结果。


\section{偏见片段定位方法}
\subsection{回复差异对齐方法}

在确认回复存在偏见风险后，本文进一步需要回答“偏见主要集中在哪些文本位置”这一问题。为此，本文对原始回复和反事实回复进行差异对齐，将整体偏见判定细化为局部风险定位。

设原始回复为 $y$，反事实回复为 $y_i^{cf}$。本文首先在句子级别对两者进行分割和对齐，确定发生显著变化的句子单元。随后，在高差异句内部进一步进行短语级或span级对齐，以定位具体变化位置。由于本文的目标是为后续重写提供最小必要修改范围，因此差异对齐不仅要找出“变了哪些句子”，还要尽量缩小到“真正承载偏见风险的局部表达”。

\subsection{候选偏见片段抽取}

基于差异对齐结果，本文从原始回复中抽取候选偏见片段集合
\begin{equation}
P(y)=\{p_1,p_2,\ldots,p_m\},
\end{equation}
其中每个片段 $p_i$ 可以是一个句子、一个短语，或一个连续的文本span。

候选片段的抽取并非简单依赖字符级差异，而是综合以下三类信号。第一，原始回复与反事实回复在该位置是否存在显著内容变化；第二，该位置是否对应显式规则或验证结果中归因为偏见来源的部分；第三，该位置是否包含明显的有害表达、刻板表述或不合理差别性建议。只有同时满足差异显著性与偏见相关性的片段，才被保留为候选偏见片段。

\subsection{片段级风险标注机制}

为了为第四章提供更细粒度的重写依据，本文进一步为每个候选偏见片段定义片段级风险强度。设片段 $p_i$ 的风险强度为 $B(p_i)$，则可表示为
\begin{equation}
B(p_i)=\eta_1 \Delta_{align}(p_i)+\eta_2 \Delta_{judge}(p_i)+\eta_3 \Delta_{harm}(p_i),
\end{equation}
其中 $\Delta_{align}(p_i)$ 表示该片段在原始回复与反事实回复之间的对齐差异强度，$\Delta_{judge}(p_i)$ 表示验证模块对该片段的偏见归因强度，$\Delta_{harm}(p_i)$ 表示该片段在有害表达维度上的风险程度，$\eta_1,\eta_2,\eta_3$ 为相应权重。

片段级风险标注的意义在于，它将样本级偏见判定进一步细化为可操作的重写对象，使得第四章无需对整条回复进行盲目改写，而能够有针对性地选择高风险区域进行局部修正。


\begin{algorithm}[htbp]
\caption{偏见片段定位与风险标注算法}
\label{alg:span_localization}
\begin{algorithmic}[1]
\Require 原始回复 $y$；反事实回复集合 $Y^{cf}(x)$；双阶段验证结果 $J$；差异归因函数 $\Delta_{align}(\cdot)$；有害性函数 $\Delta_{harm}(\cdot)$
\Ensure 候选偏见片段集合 $P(y)$ 及其风险强度 $\{B(p_i)\}$
\State 初始化候选片段集合 $P(y)\gets \emptyset$
\State 将原始回复 $y$ 切分为句子集合 $\{s_1,s_2,\ldots,s_l\}$
\For{每个句子 $s_j$}
    \State 计算 $s_j$ 与各反事实回复中对齐句子的差异强度
    \If{$s_j$ 的差异强度高于句子级阈值}
        \State 在 $s_j$ 内进行短语级或span级差异对齐
        \For{每个候选片段 $p_i$}
            \If{$p_i$ 与验证归因结果一致}
                \State 计算片段风险强度
                \[
                B(p_i)\gets \eta_1\Delta_{align}(p_i)+\eta_2\Delta_{judge}(p_i)+\eta_3\Delta_{harm}(p_i)
                \]
                \State $P(y)\gets P(y)\cup \{p_i\}$
            \EndIf
        \EndFor
    \EndIf
\EndFor
\State \Return $(P(y),\{B(p_i)\})$
\end{algorithmic}
\end{algorithm}
为将样本级偏见判定进一步转化为可用于局部校正的结构化输入，本文将片段级风险定位过程归纳为算法~\ref{alg:span_localization}。该算法在句子级和span级两层粒度上对原始回复与反事实回复进行差异对齐，并结合验证归因结果为候选片段赋予风险强度。


\section{本章小节}
本章面向黑盒大语言模型输出偏见识别问题，提出了一种由敏感属性提取、反事实prompt构造、多维差异检测、双阶段验证和片段定位组成的识别方法。该方法以输入输出行为为分析对象，通过反事实一致性分析发现潜在偏见信号，并结合规则约束与显式推理对候选偏见回复进行高置信确认。在此基础上，本文进一步输出样本级偏见标签、偏见敏感性分数和片段级风险信息，为第四章的偏见回复定向重写提供了必要的结构化输入。
